\documentclass[10pt,a4paper]{amsart}
\usepackage{amsmath,amssymb,amsthm,graphicx}
\usepackage{verbatim}
\newenvironment{reference}{\comment}{\endcomment}
%\newenvironment{reference}{}{}
\newenvironment{slogan}{\comment}{\endcomment}
\newenvironment{history}{\comment}{\endcomment}

% For commutative diagrams we use Xy-pic
\usepackage[all]{xy}

% We use 2cell for 2-commutative diagrams.
\xyoption{2cell}
% Theorem environments.
%
\theoremstyle{plain}
\newtheorem{theorem}[subsection]{Theorem}
\newtheorem{proposition}[subsection]{Proposition}
\newtheorem{lemma}[subsection]{Lemma}

\theoremstyle{definition}
\newtheorem{definition}[subsection]{Definition}
\newtheorem{example}[subsection]{Example}
\newtheorem{exercise}[subsection]{Exercise}
\newtheorem{situation}[subsection]{Situation}

\theoremstyle{remark}
\newtheorem{remark}[subsection]{Remark}
\newtheorem{remarks}[subsection]{Remarks}

\numberwithin{equation}{subsection}

% Macros
%
\def\lim{\mathop{\mathrm{lim}}\nolimits}
\def\colim{\mathop{\mathrm{colim}}\nolimits}
\def\Spec{\mathop{\mathrm{Spec}}}
\def\Hom{\mathop{\mathrm{Hom}}\nolimits}
\def\Ext{\mathop{\mathrm{Ext}}\nolimits}
\def\SheafHom{\mathop{\mathcal{H}\!\mathit{om}}\nolimits}
\def\SheafExt{\mathop{\mathcal{E}\!\mathit{xt}}\nolimits}
\def\Sch{\mathit{Sch}}
\def\Mor{\mathop{\mathrm{Mor}}\nolimits}
\def\Ob{\mathop{\mathrm{Ob}}\nolimits}
\def\Sh{\mathop{\mathit{Sh}}\nolimits}
\def\NL{\mathop{N\!L}\nolimits}
\def\CH{\mathop{\mathrm{CH}}\nolimits}
\def\proetale{{pro\text{-}\acute{e}tale}}
\def\etale{{\acute{e}tale}}
\def\QCoh{\mathit{QCoh}}
\def\Ker{\mathop{\mathrm{Ker}}}
\def\Im{\mathop{\mathrm{Im}}}
\def\Coker{\mathop{\mathrm{Coker}}}
\def\Coim{\mathop{\mathrm{Coim}}}

% Boxtimes
%
\DeclareMathSymbol{\boxtimes}{\mathbin}{AMSa}{"02}

%
% Macros for moduli stacks/spaces
%
\def\QCohstack{\mathcal{QC}\!\mathit{oh}}
\def\Cohstack{\mathcal{C}\!\mathit{oh}}
\def\Spacesstack{\mathcal{S}\!\mathit{paces}}
\def\Quotfunctor{\mathrm{Quot}}
\def\Hilbfunctor{\mathrm{Hilb}}
\def\Curvesstack{\mathcal{C}\!\mathit{urves}}
\def\Polarizedstack{\mathcal{P}\!\mathit{olarized}}
\def\Complexesstack{\mathcal{C}\!\mathit{omplexes}}
% \Pic is the operator that assigns to X its picard group, usage \Pic(X)
% \Picardstack_{X/B} denotes the Picard stack of X over B
% \Picardfunctor_{X/B} denotes the Picard functor of X over B
\def\Pic{\mathop{\mathrm{Pic}}\nolimits}
\def\Picardstack{\mathcal{P}\!\mathit{ic}}
\def\Picardfunctor{\mathrm{Pic}}
\def\Deformationcategory{\mathcal{D}\!\mathit{ef}}
\usepackage{cleveref}
\crefname{originalreference}{original reference}{original references}
\Crefname{originalreference}{Original reference}{Original references}
\makeatletter
\newlabel{definition-rig-smooth-homomorphism}{{0GAI}{source}{}{sourcepointer-1}{}}
\newlabel{definition-rig-smooth-homomorphism@cref}{{[originalreference][0][]0GAI}{[1][0][]source}{}{}{}}
\makeatother
\crefname{originalequation}{original equation}{original equations}
\Crefname{originalequation}{Original equation}{Original equations}
\makeatletter
\newlabel{equation-C-prime}{{0AL4}{source}{}{sourcepointer-2}{}}
\newlabel{equation-C-prime@cref}{{[originalequation][0][]0AL4}{[1][0][]source}{}{}{}}
\makeatother
\crefname{originalreference}{original reference}{original references}
\Crefname{originalreference}{Original reference}{Original references}
\makeatletter
\newlabel{lemma-zero-ext-1-after-modding-out}{{0GAL}{source}{}{sourcepointer-3}{}}
\newlabel{lemma-zero-ext-1-after-modding-out@cref}{{[originalreference][0][]0GAL}{[1][0][]source}{}{}{}}
\makeatother
\crefname{originalreference}{original reference}{original references}
\Crefname{originalreference}{Original reference}{Original references}
\makeatletter
\newlabel{remark-base-change}{{06FQ}{source}{}{sourcepointer-4}{}}
\newlabel{remark-base-change@cref}{{[originalreference][0][]06FQ}{[1][0][]source}{}{}{}}
\makeatother
\title{Rig-smooth homomorphisms are stable under base change}
\author{The Stacks Project Authors}
\date{stacks-project:restricted.tex:local-snapshot-20261012; GNU FDL 1.2 or later}
\begin{document}
\maketitle
\noindent\textit{Editorial note.} All delivered mathematical text is the author's own, copied from the pinned Stacks Project chapter file restricted.tex; no statement or proof step was rewritten, reordered or removed.\par\medskip
\noindent\textit{Editorial note.} The counted proof is the author's complete argument: the counted statement is the author lemma environment labelled lemma-base-change-rig-smooth-homomorphism; the proof environment immediately following it is its own proof\par\medskip
\noindent\textit{Editorial note.} The label lemma-base-change-rig-smooth-homomorphism names the counted statement, and the proof environment immediately following it in the source is delivered verbatim as its proof.\par\medskip
\noindent\textit{Editorial note.} The counted statement is carried by the author's own lemma environment in the chapter source.\par\medskip
\noindent\textit{Allowed standard prerequisites (editorial).}\begin{itemize}
\item the base change setup for pairs (A1, I1) to (A2, I2) of Noetherian rings with V(I1 A2) = V(I2) and the base change algebra B2 (remark-base-change, equation-C-prime)
\item the Ext\textasciicircum{}1 annihilation criterion and its persistence under base change (definition-rig-smooth-homomorphism and lemma-zero-ext-1-after-modding-out, delivered as its own card)
\item commensurability of ideals with equal radical in a Noetherian ring: I2\textasciicircum{}c is contained in I1 A2 for some c >= 0
\end{itemize}\medskip
\ifdefined\hypertarget\hypertarget{sourcepointer-1}{}\fi
\noindent\textit{Source reference.} Original reference 0GAI: sources/stacks-restricted/pinned-original-source.tex, source lines 831--831. Prerequisite pointer invoked at 1055; it names a standard fact of the same chapter that the counted argument uses, and every step of the counted proof built on it is retained verbatim in the delivered body.\par\medskip
\ifdefined\hypertarget\hypertarget{sourcepointer-2}{}\fi
\noindent\textit{Source reference.} Original equation 0AL4: sources/stacks-restricted/pinned-original-source.tex, source lines 326--326. Prerequisite pointer invoked at 1046; it names a standard fact of the same chapter that the counted argument uses, and every step of the counted proof built on it is retained verbatim in the delivered body.\par\medskip
\ifdefined\hypertarget\hypertarget{sourcepointer-3}{}\fi
\noindent\textit{Source reference.} Original reference 0GAL: sources/stacks-restricted/pinned-original-source.tex, source lines 1016--1016. Prerequisite pointer invoked at 1054; it names a standard fact of the same chapter that the counted argument uses, and every step of the counted proof built on it is retained verbatim in the delivered body.\par\medskip
\ifdefined\hypertarget\hypertarget{sourcepointer-4}{}\fi
\noindent\textit{Source reference.} Original reference 06FQ: sources/stacks-restricted/pinned-original-source.tex, source lines 426--426. Prerequisite pointer invoked at 1048; it names a standard fact of the same chapter that the counted argument uses, and every step of the counted proof built on it is retained verbatim in the delivered body.\par\medskip
\begin{lemma}
\label{lemma-base-change-rig-smooth-homomorphism}
Let $A_1 \to A_2$ be a map of Noetherian rings. Let $I_i \subset A_i$
be an ideal such that $V(I_1A_2) = V(I_2)$. Let $B_1$ be in
(\ref{equation-C-prime}) for $(A_1, I_1)$.
Let $B_2$ be the base change of $B_1$ as in
Remark \ref{remark-base-change}.
If $B_1$ is rig-smooth over $(A_1, I_1)$,
then $B_2$ is rig-smooth over $(A_2, I_2)$.
\end{lemma}

\begin{proof}
Follows from Lemma \ref{lemma-zero-ext-1-after-modding-out} and
Definition \ref{definition-rig-smooth-homomorphism}
and the fact that $I_2^c$ is contained in $I_1A_2$ for some $c \geq 0$
as $A_2$ is Noetherian.
\end{proof}

\end{document}
